
\textbf{SKYFILMS}, una empresa colombiana líder en producción audiovisual con drones, se compromete a brindar servicios de alta calidad de manera segura y responsable. Con el objetivo de garantizar la seguridad de nuestras operaciones, hemos desarrollado este Manual del Sistema de Gestión de Seguridad Operacional (SMS), alineado con los requisitos establecidos en el Reglamento Aeronáutico Civil (RAC) en sus versiones 100 y 219.\\

Este manual establece las directrices para la gestión sistemática de la seguridad operacional, documentando las estructuras organizacionales, líneas de responsabilidad, políticas y procedimientos necesarios para mantener nuestras operaciones bajo control. Para \textbf{SKYFILMS}, el SMS no es un componente aislado, sino un enfoque de negocio fundamental que se integra en todos los procesos de planificación, ejecución y evaluación de nuestras misiones audiovisuales.\\

A través de este documento, la organización consolida los cuatro pilares fundamentales del sistema: la Política y Objetivos de Seguridad Operacional, la Gestión de Riesgos, el Aseguramiento de la Seguridad y la Promoción de la Seguridad Operacional. Nuestro propósito es transitar de un cumplimiento normativo básico hacia una gestión proactiva y predictiva, garantizando la identificación temprana de peligros y la implementación de defensas eficaces que mantengan los riesgos en un nivel tan bajo como sea razonablemente práctico.


\section{Alcance}

El presente manual establece el alcance del Sistema de Gestión de Seguridad Operacional (SMS) de \textbf{SKYFILMS}, el cual es de estricto cumplimiento para todo el personal y aplicable a todas las actividades relacionadas con la prestación de nuestros servicios aéreos.\\

\textbf{Alcance Operacional:}

El SMS aplica a todas las operaciones aéreas realizadas por \textbf{SKYFILMS} utilizando Sistemas de Aeronaves No Tripuladas (UAS). Esto incluye las fases de planificación, preparación, ejecución y post-vuelo. Nuestro ámbito de operación abarca entornos diversos y complejos, incluyendo zonas urbanas, áreas rurales y operaciones sobre multitudes, empleando una variedad de aeronaves no tripuladas según las autorizaciones correspondientes.\\

\textbf{Alcance Organizacional:}\\
La seguridad operacional es una responsabilidad transversal en \textbf{SKYFILMS}. Por lo tanto, los lineamientos, políticas y procedimientos descritos en este manual aplican a todos los niveles estructurales de la organización, vinculando expresamente a:\\

\begin{itemize}
    \item \textbf{Gerencia General} (como Ejecutivo Responsable de la provisión de recursos y del SMS).

    \item \textbf{Gerente de Seguridad Operacional} (responsable de la administración del sistema).

    \item Jefe de Pilotos.

    \item Pilotos al mando, observadores y cualquier personal técnico o de apoyo involucrado en la operación.
\end{itemize}


\textbf{Aplicabilidad:}\\
Para garantizar un enfoque sistémico, los componentes del SMS se integran de manera obligatoria en los siguientes procesos de la organización:\\

\begin{itemize}
    \item \textbf{Operaciones:} Planificación, supervisión, autorización y ejecución segura de todos los vuelos UAS.

    \item \textbf{Mantenimiento:} Procesos de inspección, verificación pre-vuelo, conservación y gestión de la aeronavegabilidad continua de los equipos y estaciones de control.

    \item \textbf{Gestión de Riesgos:} Identificación proactiva y predictiva de peligros, evaluación de la probabilidad y severidad, y aplicación de medidas de mitigación.

    \item \textbf{Aseguramiento de la Seguridad:} Reporte voluntario y obligatorio de eventos de seguridad, investigación, análisis de tendencias y auditorías internas para la mejora continua.

    \item \textbf{Promoción de la Seguridad:} Capacitación continua del personal, evaluación de competencias y promoción activa de la comunicación y una cultura justa.
\end{itemize}


\section{Objetivo del Manual}
El objetivo principal de este manual es establecer un marco de trabajo que permita identificar, evaluar y mitigar los riesgos asociados a nuestras operaciones con drones. Buscamos garantizar el cumplimiento de la normativa vigente, prevenir incidentes y accidentes, proteger a las personas y bienes, y promover una cultura de seguridad en toda la organización.

